\documentclass{anafw}

% ─── Metadados do documento ────────────────────────────────────────────────────
\consultoria{Nome da Consultoria}
\logoConsultoria{logo.png}
\documentoId{ANA-FW-001}
\tituloAnalise{Análise de Firmware}
\subtituloAnalise{[Nome do produto / subsistema analisado]}
\projetoNome{[Nome do Projeto]}
\clienteNome{[Nome do Cliente]}
\responsavel{[Engenheiro Responsável]}
\revisao{Rev A}
\plataforma{[MCU / Framework — ex.: ESP32 / Arduino]}
\firmwareAnalisado{[Versão / commit analisado]}
\repositorio{[Repositório / caminho]}

% ─── Assinaturas ───────────────────────────────────────────────────────────────
\assinaturaResponsavel{[Nome do Responsável]}{[Cargo]}
\assinaturaCliente{[Nome do Representante do Cliente]}{[Cargo] \textbar\ [Empresa do Cliente]}

\begin{document}

\fazerCapa


%% ════════════════════════════════════════════════════════════════
%%  RESUMO EXECUTIVO
%% ════════════════════════════════════════════════════════════════
\begin{resumobox}
O firmware analisado \textbf{[cumpre / cumpre parcialmente]} sua função atual de
\textbf{[descrever em 1 frase o que o dispositivo faz]}.
Foram identificados \textbf{[N]} pontos de melhoria, sendo \textbf{[N]} críticos
(corrigem bugs ou riscos reais), sem alterar o comportamento desejado do produto.
A evolução proposta de maior valor é \textbf{[descrever]}.
\end{resumobox}


%% ════════════════════════════════════════════════════════════════
%%  1. VISÃO GERAL DO SISTEMA
%% ════════════════════════════════════════════════════════════════
\section{1. Visão Geral do Sistema}

\textbf{[Descrever em 1--2 parágrafos o propósito do dispositivo: o que controla,
quais atuadores e sensores possui, e como o usuário interage com ele.]}

\subsection{Periféricos e pinos}

\begin{listaItens}
  \item \textbf{[Atuador 1]:} [função no sistema].
  \item \textbf{[Sensor / entrada]:} [função no sistema].
  \item \textbf{[Indicador / saída]:} [função no sistema].
  \item \textbf{[Interface de comando — ex.: RF, botão, UART]:} [função no sistema].
\end{listaItens}

\subsection{Arquivos / módulos}

\begin{tabularx}{\linewidth}{@{} l X @{}}
  \cabArquivos
  \texttt{[main.\,c/.ino]} & Ponto de entrada: \texttt{setup()} e \texttt{loop()}. \\[3pt]
  \texttt{[modulo\_a.h]}   & [Responsabilidade do módulo]. \\[3pt]
  \texttt{[modulo\_b.h]}   & [Responsabilidade do módulo]. \\[3pt]
  \texttt{[config.h]}      & [Constantes, pinos e parâmetros]. \\
  \bottomrule
\end{tabularx}


%% ════════════════════════════════════════════════════════════════
%%  2. O QUE O FIRMWARE FAZ HOJE
%% ════════════════════════════════════════════════════════════════
\section{2. O que o Firmware Faz Hoje}

\subsection{Inicialização — \texttt{setup()}}

\begin{listaItens}
  \item {[Passo de inicialização, ex.: configura Serial / pinos / periféricos]}.
  \item {[Passo de inicialização]}.
\end{listaItens}

\subsection{Laço principal — \texttt{loop()}}

\begin{listaItens}
  \item {[Ação executada a cada iteração]}.
  \item {[Tratamento de evento / comando recebido e o que ele dispara]}.
\end{listaItens}

\subsection{Estados internos}

\begin{tabularx}{\linewidth}{@{} l X @{}}
  \toprule
  \rowcolor{primario}
  \textcolor{white}{\textbf{Estado}} & \textcolor{white}{\textbf{Significado}} \\
  \midrule
  \texttt{[ESTADO\_A]} & [Descrição]. \\[3pt]
  \texttt{[ESTADO\_B]} & [Descrição]. \\
  \bottomrule
\end{tabularx}


%% ════════════════════════════════════════════════════════════════
%%  3. DIAGRAMA DO FLUXO ATUAL
%% ════════════════════════════════════════════════════════════════
\section{3. Diagrama do Fluxo Atual}

\begin{codigo}
            +----------- loop() -----------+
            | le entradas / eventos        |
            | comando disponivel?          |
            +--------------+---------------+
                           | sim
              +------------+------------+
              |            |            |
          [comando A]  [comando B]  [comando C]
              |            |            |
          [acao A]     [acao B]     (reservado)
\end{codigo}

\noindent\textbf{[Observação:]} [apontar, se for o caso, que o fluxo é bloqueante,
que eventos são ignorados durante operações longas, etc.]


%% ════════════════════════════════════════════════════════════════
%%  4. PONTOS DE MELHORIA
%% ════════════════════════════════════════════════════════════════
\section{4. Pontos de Melhoria}

Ordenados por impacto. Nenhum item desta seção altera o comportamento desejado do
produto --- apenas corrigem defeitos ou melhoram a qualidade do código.

\subsection{\nivelCritico\ Corrigem bugs ou riscos reais}

\begin{listaItens}
  \item \textbf{[Título do problema].} [Descrição do defeito e por que é crítico.]
        \\ \textit{Ação:} [correção recomendada].
  \item \textbf{[Título do problema].} [Descrição.]
        \\ \textit{Ação:} [correção recomendada].
\end{listaItens}

\subsection{\nivelEstrutural\ Legibilidade e manutenção}

\begin{listaItens}
  \item \textbf{[Título].} [Ex.: números mágicos, nomes enganosos, falta de \texttt{enum},
        ausência de \texttt{include guards}.] \\ \textit{Ação:} [recomendação].
  \item \textbf{[Título].} [Descrição.] \\ \textit{Ação:} [recomendação].
\end{listaItens}

\subsection{\nivelRobustez\ Resiliência}

\begin{listaItens}
  \item \textbf{[Título].} [Ex.: código bloqueante, ausência de debounce / timeout / watchdog.]
        \\ \textit{Ação:} [recomendação].
\end{listaItens}

\subsection{Resumo priorizado}

\begin{tabularx}{\linewidth}{@{} c X c c @{}}
  \cabMelhorias
  1 & [Ponto de melhoria]                    & \nivelCritico    & \esfBaixo \\[3pt]
  2 & [Ponto de melhoria]                    & \nivelCritico    & \esfBaixo \\[3pt]
  3 & [Ponto de melhoria]                    & \nivelEstrutural & \esfMedio \\[3pt]
  4 & [Ponto de melhoria]                    & \nivelEstrutural & \esfBaixo \\[3pt]
  5 & [Ponto de melhoria]                    & \nivelRobustez   & \esfAlto  \\
  \bottomrule
\end{tabularx}

\legendaCriticidade


%% ════════════════════════════════════════════════════════════════
%%  5. PLANO DE IMPLEMENTAÇÃO POR ETAPAS
%% ════════════════════════════════════════════════════════════════
\section{5. Plano de Implementação por Etapas}

\begin{tabularx}{\linewidth}{@{} l X X l @{}}
  \cabEtapas
  Etapa 1 & Refactor seguro (mesmo comportamento) & Itens críticos e estruturais  & ---       \\[3pt]
  Etapa 2 & Correções de hardware / pinos         & [Itens]                       & Etapa 1   \\[3pt]
  Etapa 3 & [Nova funcionalidade — núcleo]        & [Itens]                       & Etapa 1   \\[3pt]
  Etapa 4 & [Robustez — máquina de estados]       & [Itens]                       & Etapas 1--3 \\
  \bottomrule
\end{tabularx}


%% ════════════════════════════════════════════════════════════════
%%  6. EVOLUÇÃO PROPOSTA (DESCRIÇÃO)
%% ════════════════════════════════════════════════════════════════
\section{6. Evolução Proposta}

\subsection{Requisito}

\textbf{[Descrever a nova funcionalidade desejada, do ponto de vista do usuário /
instalador. O que muda em relação ao comportamento atual.]}

\begin{premissabox}
[Apontar a premissa que precisa de confirmação do cliente e como ela muda a
modelagem. Ex.: definição exata de um ``par'', limites de capacidade, etc.]
\end{premissabox}

\subsection{Abordagem recomendada}

\begin{listaItens}
  \item \textbf{Persistência:} [ex.: usar \texttt{Preferences}/NVS no ESP32 em vez de emular EEPROM].
  \item \textbf{Modelo de dados:} [estrutura e capacidade].
  \item \textbf{Máquina de estados:} [estados do novo fluxo e gatilhos].
  \item \textbf{Feedback ao usuário / timeout / migração:} [como sinalizar e como não quebrar instalações existentes].
\end{listaItens}

\subsection{Esboço de modelo de dados}

\begin{codigo}[language=C]
#define MAX_ITENS 2

typedef struct {
    bool     valido;     // slot gravado?
    uint32_t cod_a;      // [funcao principal]
    uint32_t cod_b;      // [funcao secundaria]
} item_t;

item_t itens[MAX_ITENS];
\end{codigo}


%% ════════════════════════════════════════════════════════════════
%%  7. ASSINATURAS
%% ════════════════════════════════════════════════════════════════
\section{7. Emissão e Aprovação}

\fazerAssinaturas

\end{document}
